% 原书 p.214–p.216（PDF p.222–p.224）
% 第 8 章 8.4 节：量子上同调群（Quantum Cohomology Group）
% 说明：本节写到 p.224（8.4 节结束）；8.5–8.7 由另一译者续写（ch08-s05 起）。

\section{量子上同调群（Quantum Cohomology Group）}
% >>> p222

\begin{keytakeaways}
\begin{itemize}
  \item 几何图像：量子上同调群 $\mathrm{QH}^*(M,\omega)$ 是经典上同调群与 Novikov 域的张量积，
        它的乘法（量子乘积 $*$）在经典杯积之外多出由 $J$-全纯曲线计数给出的「形变项」。
  \item 数学工具：结构常数 $c_{ij}^k$（经典）与 $a_{ij}^k\in\Lambda^{\mathbf{k},\Gamma_\omega}$
        （量子）；分次（$\deg(e_1)=2$、$\deg(T)=2n+2$）；Poincaré 对偶配对
        $\langle x,y\rangle=\int_{[M]}x\cup y$ 及其在 $\mathrm{QH}^*$ 上的推广 (8.35)。
  \item 本节结论：$\mathrm{QH}^*(\CP^n,\omega_{\mathrm{FS}})=\mathbf{k}[s,T]/\langle s^{n+1}=T\rangle$，
        $\mathrm{QH}^*(S^2\times S^2)$ 中 $a*a=T^{[\mathrm{pt}]\times S^2}\cdot1$、
        $b*b=T^{S^2\times[\mathrm{pt}]}\cdot1$（用增强 Novikov 环才能区分二者）；
        $(\mathrm{QH}^*,*,\langle\cdot,\cdot\rangle)$ 是结合 Frobenius 代数（事实 8.4.1）。
  \item 阅读要点：量子乘积把「经典部分」与「形变部分」分开写（(8.30)），形变部分才是曲线的
        计数信息；$\mathrm{QH}^*$ 中更多元素变得可逆（注 8.4.1），这是它与经典上同调环的本质差别。
\end{itemize}
\end{keytakeaways}

让我们从下述定义开始。

\begin{wdefn}{8.4.1}
  给定辛流形 $(M,\omega)$，其\textbf{量子上同调群}（quantum cohomology groups）定义为
  \[
    \mathrm{QH}^*(M,\omega):=H^*(M;\mathbf{k})\otimes_{\mathbf{k}}
      \Lambda^{\mathbf{k},\Gamma_\omega},
  \]
  其中 $\Lambda^{\mathbf{k},\Gamma_\omega}$ 是 8.1 节定义的 Novikov 环。
\end{wdefn}
% [Tutor Note: 易错预警（分次与维数）——(1) $\mathrm{QH}^*(M,\omega)$ 是 Novikov 域
%  $\Lambda^{\mathbf{k},\Gamma_\omega}$ 上的有限维向量空间，其 $\Lambda$-秩等于
%  $\sum_*\dim_{\mathbf{k}}H^*(M;\mathbf{k})$，但作为 $\mathbf{k}$-向量空间它是无穷维的
%  （每个经典类都可以乘上任意 $T^\lambda$）。(2) 分次结构需要约定形式变量的度数：
%  按 8.1 节 (8.28) 前的约定取 $\deg(T^A):=2c_1(TM)(A)$，于是 $\deg(T^\lambda)=2\lambda$
%  只在单调情形由 $\lambda\mapsto c_1$ 良定义；对一般的 $\omega$，$\mathrm{QH}^*$ 只有由 $\deg$ 给出的
%  「偶次」分次（见 (8.32) 下方 $\deg(e_1)=2$、$\deg(T)=2n+2$ 的例子）。逐符号核对分次时
%  请以所用的 $\deg(T)$ 约定为准，这是文献中最易出错的地方。]

显然，$\mathrm{QH}^*(M,\omega)$ 是 Novikov 域 $\Lambda^{\mathbf{k},\Gamma_\omega}$ 上的有限维向量空间。
我们的目标是用 $\mathrm{QH}^*(M,\omega)$ 研究辛流形 $(M,\omega)$ 的性质。由 8.2 节的非阿基米德线性代数，
$\mathrm{QH}^*(M,\omega)$ 在 $\Lambda^{\mathbf{k},\Gamma_\omega}$ 上的线性代数已经清楚。经典上同调群强有力
的杯积在 $\mathrm{QH}^*(M,\omega)$ 上也有一个类比。更确切地说，$\mathrm{QH}^*(M,\omega)$ 上的杯积以如下方式
加强 $H^*(M;\mathbf{k})$ 上的经典杯积。

设 $H^*(M;\mathbf{k})=\operatorname{span}_{\mathbf{k}}\langle e_1,\cdots,e_n\rangle$，其中
$n=\sum_*\dim_{\mathbf{k}}H^*(M;\mathbf{k})$。对任何基元素 $e_i$ 与 $e_j$，经典杯积计算如下：
\begin{equation}\tag{8.29}
  e_i\cup e_j=\sum_{k=1}^n c_{ij}^k e_k
    \qquad\text{其中 }c_{ij}^k\in\mathbf{k},
\end{equation}
其中 $c_{ij}^k$ 表示杯积的结构常数。现在，把 $\{e_1,\cdots,e_n\}$ 看作 $\mathrm{QH}^*(M,\omega)$ 在
$\Lambda^{\mathbf{k},\Gamma_\omega}$ 上的一组基，则\textbf{量子乘积}（quantum product，其更精确的
定义将在稍后的 8.7 节给出，记作 $*$）如下计算：
\begin{equation}\tag{8.30}
  e_i*e_j=\sum_{k=1}^n a_{ij}^k e_k
    =e_i\cup e_j+(\text{形变部分}),
\end{equation}
其中 $a_{ij}^k\in\Lambda^{\mathbf{k},\Gamma_\omega}$ 是量子乘积的结构常数，第二个等号来自对各项的
重新整理。换言之，与 (8.29) 相比，量子乘积给出了额外的「形变」项。尽管如此，装备了 $*$ 之后，
量子上同调群 $\mathrm{QH}^*(M,\omega)$ 仍是一个代数，其单位为 $1\in\mathbf{k}\,(=H^0(M;\mathbf{k}))$。
% [Tutor Note: 易错预警——(8.30) 中的 $a_{ij}^k$ 一般落在 $\Lambda^{\mathbf{k},\Gamma_\omega}$
%  中（不只是在 $\mathbf{k}$ 中），其 $T^{>0}$ 部分正是「形变项」，由计数经过 $e_i,e_j$ 对应
%  的 Poincaré 对偶类的 $3$ 点 Gromov--Witten 不变量给出；这就是为什么量子乘积的严格定义要
%  等到 8.7 节（先有 8.5 节的 GW 不变量）。此外 (8.30) 把「经典部分」与「形变部分」相加，隐含
%  地使用了 $H^*(M;\mathbf{k})\hookrightarrow\mathrm{QH}^*(M,\omega)$（$T^0$ 系数）这一嵌入，切勿把
%  $e_i\cup e_j$ 与 $e_i*e_j$ 混淆。]

在辛几何中，$\mathrm{QH}^*(M,\omega)$ 的代数性质有助于发现拓扑刚性（本节末将给出一个应用）。

在给出 $\mathrm{QH}^*(M,\omega)$ 的严格定义之前，先看两个例子，以展示量子上同调群的计算结果。

\begin{wexam}{8.4.1}
  考虑 $(M,\omega)=(\CP^n,\omega_{\mathrm{FS}})$。众所周知，其上同调群
  $H^*(\CP^n;\mathbf{k})=\operatorname{span}_{\mathbf{k}}\langle e_0,e_1,\cdots,e_n\rangle$，
  其中 $e_0=1$，$e_1=\mathrm{PD}([L])$ 是 $\CP^n$ 中标准直线所代表的类的 Poincaré 对偶，而对
  $i$ 次有 $e_i=e_1^i=e_1\cup\cdots\cup e_1$。事实上，作为一个代数，
  $H^*(\CP^n;\mathbf{k})=\mathbf{k}[s]/\langle s^{n+1}=0\rangle$（多项式环的商）。对量子上同调
  群，我们有
  \begin{equation}\tag{8.31}
    \mathrm{QH}^*(\CP^n,\omega_{\mathrm{FS}})
      =\frac{\mathbf{k}[s,T]}{\langle s^{n+1}=T\rangle},
  \end{equation}
  其中 $T$ 是一个新的形式变量。蕴涵 (8.31) 的关键杯积关系如下：
  \begin{equation}\tag{8.32}
    \underbrace{e_1*\cdots*e_1}_{(n+1)\text{ 个}}
      =\underbrace{e_1\cup\cdots\cup e_1}_{(n+1)\text{ 个}}
      +\underbrace{T\cdot e_0}_{\text{形变部分}}
      =T\cdot e_0.
  \end{equation}
  也值得一提的是 (8.32) 中出现的次数。由于 $\deg(e_1)=2$，我们有 $\deg(T)=2n+2$。\hfill$\square$
\end{wexam}
% [Tutor Note: 易错预警——(8.32) 中 $e_1^{\cup(n+1)}=0$（因 $H^{2(n+1)}(\CP^n)=0$），故右端只剩
%  $T\cdot e_0$；这正是 (8.31) 中关系 $s^{n+1}=T$ 的来源。$\deg(T)=2n+2$ 与
%  $\deg(e_1^{n+1})=2n+2$ 一致，是分次结构自洽的检验点；若用 $c_1$ 的约定，
%  $c_1(T\CP^n)=n+1$、$L$ 为直线类，则 $\deg(T)=2c_1(L)=2n+2$，两者相符。]

% >>> p223
\begin{wexam}{8.4.2}
  考虑 $(M,\omega)=(S^2\times S^2,\omega_{\mathrm{prod}})$。其上同调群有一组基
  $\{1,a,b,a\cup b\}$，其中 $a$ 与 $b$ 分别是由两个 $S^2$-因子代表的上同调类。由例 8.4.1
  （应用于 $\CP^1$），我们有
  \begin{equation}\tag{8.33}
    a\cup a=b\cup b=0\qquad\text{但}\qquad a*a=b*b=T\cdot1.
  \end{equation}
  有趣的是，不同版本的系数环或系数域可以帮助我们区分 $a*a$ 与 $b*b$，尽管它们在 (8.33) 中的
  计算结果相同。确实，把 $\Lambda^{\mathbf{k},\Gamma_\omega}$ 加强为 (8.8) 的 Novikov 条件下由
  (8.6) 中元素所定义的 $\Lambda^{\mathbf{k},c}$，并考虑
  $H^*(M;\mathbf{k})\otimes_{\mathbf{k}}\Lambda^{\mathbf{k},c}$。上述同样的计算表明
  \[
    a*a=T^{[\mathrm{pt}]\times S^2}\cdot1
    \qquad\text{且}\qquad
    b*b=T^{S^2\times[\mathrm{pt}]}\cdot1.
  \]
  因此，在这个增强的 Novikov 环上，$a*a$ 与 $b*b$ 是不同的。\hfill$\square$
\end{wexam}
% [Tutor Note: 易错预警——这里 $T^{[\mathrm{pt}]\times S^2}$ 与 $T^{S^2\times[\mathrm{pt}]}$
%  的指数是 $H_2^S(S^2\times S^2)$ 中的**同调类**而不是标量，这正是 8.1 节修改 (ii) 的
%  $\Lambda^{\mathbf{k},c}$ 相对于 $\Lambda^{\mathbf{k},\Gamma_\omega}$ 的用处：在
%  $\Lambda^{\mathbf{k},\Gamma_\omega}$ 中两者的指数都等于 $\omega(\cdot)=1$，无法区分；
%  只有把 $\ker[\omega]$ 方向的类（这里 $[S^2]\times[\mathrm{pt}]$ 与
%  $[\mathrm{pt}]\times[S^2]$ 之差）也记入变量，才能区分 $a*a$ 与 $b*b$。这与 8.1 节
%  (8.7) 中 $S^B$ 的构造完全一致，核对时请把 $T^\lambda$ 与 $S^B$ 两个因子对应起来。]

\begin{wrem}{8.4.1}
  这里给出经典上同调群与量子上同调群之间的一个本质差别：$\mathrm{QH}^*(M,\omega)$ 中有更多元素变得
  可逆。例如，在例 8.4.1 中，基元素 $e_1$ 在 $H^*(\CP^n;\mathbf{k})$ 中\textbf{不}可逆，但在
  $\mathrm{QH}^*(\CP^n,\omega_{\mathrm{FS}})$ 中可逆，且
  $e_1^{-1}=T^{-1}\cdot e_1^n=T^{-1}\cdot e_n$。更一般地，由于
  $s^{n+1}-T\in\mathbf{k}[s,T]$ 是不可约多项式，由 (8.32)，量子上同调
  $\mathrm{QH}^*(\CP^n,\omega_{\mathrm{FS}})$ 作为某个极大理想的商是一个域，从而其中每个元素都可逆。
\end{wrem}

回忆在紧定向流形 $M$ 上，经典上同调群 $H^*(M;\mathbf{k})$ 上有一个配对结构：对任何满足
$\deg(x)+\deg(y)=\dim M$ 的 $x,y\in H^*(M;\mathbf{k})$，定义
\begin{equation}\tag{8.34}
  \langle x,y\rangle:=\int_{[M]}x\cup y\ \bigl(=x(\mathrm{PD}(y))\bigr)\in\mathbf{k},
\end{equation}
其中第二个等号来自 Poincaré 对偶的定义。在这一框架下，$\mathrm{QH}^*(M,\omega)$ 上也有一个配对
（其取值落在 $\mathbf{k}$ 中），定义如下。

% >>> p224
注意由线性性，只需对 $\mathrm{QH}^*(M,\omega)$ 中形如 $\lambda x$、$\eta y$ 的单项式元素定义配对，
其中 $\lambda,\eta\in\Lambda^{\mathbf{k},c}$，$x,y\in H^*(M;\mathbf{k})$ 且
$\deg(x)+\deg(y)=\dim M$。于是定义
\begin{equation}\tag{8.35}
  \langle\lambda x,\eta y\rangle
    :=(\lambda\eta)_0\int_{[M]}x\cup y\in\mathbf{k},
\end{equation}
其中 $(\lambda\eta)_0$ 表示取 $T^0$-项的系数（换言之，常数项）。下面是一个例子。
% [Tutor Note: 易错预警——(1) $(\lambda\eta)_0$ 只取 $\lambda\eta$ 的 $T^0$ 系数：由于
%  $\Lambda^{\mathbf{k},c}$ 的元素允许无穷多项，若不做这一步投影，配对可能不收敛/无定义；
%  (8.8) 的 Novikov 条件保证每个指数的系数有限，故取常数项是良定义的。(2) 这个配对对两个
%  变元都是 $\Lambda^{\mathbf{k},c}$-线性的（不是共轭线性），这一点与 Hermite 内积不同，
%  也是 Frobenius 代数公理 $\langle a*b,c\rangle=\langle a,b*c\rangle$ 得以成立的原因。]

\begin{wexam}{8.4.3}
  对 $(\CP^n,\omega_{\mathrm{FS}})$，基元素 $e_i$（$0\leqslant i\leqslant n$）满足
  \[
    \langle e_i,e_{n-i}\rangle=1,
  \]
  而当 $j\ne n-i$ 时 $\langle e_i,e_j\rangle=0$。注意在 (8.35) 中，唯一非平凡的部分是
  「系数配对」$\langle\lambda e_i,\eta e_{n-i}\rangle\mapsto(\lambda\eta)_0$
  （其中 $x=e_i$、$y=e_{n-i}$，$0\leqslant i\leqslant n$）。由本章引言的讨论，由于
  $(\CP^n,\omega_{\mathrm{FS}})$ 是单调的，相应的 Novikov 环同构于完备化
  $\mathbf{k}[\![T]\!]$；于是容易验证配对 $\langle\lambda,\eta\rangle\mapsto(\lambda\eta)_0$
  是非退化的——即若 $\langle\lambda,\eta\rangle=0$，则 $\lambda,\eta$ 中必有一个为零。这蕴涵
  (8.35) 中的配对也是非退化的。\hfill$\square$
\end{wexam}

下面是本节的主要结果，它总结了量子上同调群的代数性质。

\whead{事实}{8.4.1}（[154] 中的命题 11.11）\quad 设 $(M,\omega)$ 是紧辛流形，则其量子上同调群
$\mathrm{QH}^*(M,\omega)$ 与量子乘积 $*$、以及 (8.35) 中定义的配对 $\langle\cdot,\cdot\rangle$ 一起，
构成一个结合的\textbf{Frobenius 代数}（Frobenius algebra）：即 $\langle\cdot,\cdot\rangle$ 是
非退化双线性形式，且满足
\[
  \langle a*b,c\rangle=\langle a,b*c\rangle.
\]
% [Tutor Note: 术语依据与含糊点拨——(1) 「Frobenius 代数」是有限维结合代数配以非退化结合双线性
%  形式（$\langle a*b,c\rangle=\langle a,b*c\rangle$）的结构；量子上同调是它的标准例子。(2)
%  结合性 $(\mathrm{QH}^*,*)$ 的结合律等价于 Gromov--Witten 不变量的 WDVV 方程（Witten--Dijkgraaf--
%  Verlinde--Verlinde），这正是 8.7 节的主题；本节只陈述 Frobenius 代数这一「输出」，它的证明
%  依赖 8.5 节的 GW 不变量定义与 8.6–8.7 节的粘合公式。阅读 WDVV 前务必确认三项前置约定：
%  (a) 量子乘积的结构常数取哪一组基；(b) $\deg(T)$ 的约定（本译本沿用 $\deg(T^A)=2c_1(TM)(A)$）；
%  (c) 配对 (8.35) 中「取 $T^0$ 系数」的投影 $(\cdot)_0$。三者任一改变，WDVV 方程的写法都会
%  随之改变（可能相差符号或因子 $2$）。]

\begin{wrem}{8.4.2}
  由于 $\mathrm{QH}^*(M,\omega)$ 上的配对 $\langle\cdot,\cdot\rangle$ 非退化，可以如下定义元素
  $x\in\mathrm{QH}^*(M,\omega)$ 的\textbf{对偶}：唯一的元素 $y\in\mathrm{QH}^*(M,\omega)$ 使得
  $\langle x,y\rangle=1$。
\end{wrem}

为结束本节，让我们说明 $\mathrm{QH}^*(M,\omega)$ 上的代数性质如何蕴涵辛流形 $(M,\omega)$ 自身的一些
特殊性质。

\begin{wexer}{8.4.1}（[78] 第 2.3 节）
  域 $\mathbf{k}$ 上的有限维代数 $\mathcal{A}$ 若可分解为若干域的直和，则称为\textbf{半单的}
  （semi-simple）。由 Abram 判据（见 [9]），对 Frobenius 代数 $\mathcal{A}$，它半单当且仅当
  \[
    \mathcal{E}:=\sum_{\text{基 }e_i}e_i*e_i^{\#}
  \]
  在 $\mathcal{A}$ 中可逆，其中 $e_i^{\#}$ 是 $e_i$ 关于 $\mathcal{A}$ 中配对
  $\langle\cdot,\cdot\rangle$ 的对偶。证明：对 $(\CP^n,\omega_{\mathrm{FS}})$ 与
  $(S^2\times S^2,\omega_{\mathrm{prod}})$ 两者，其量子上同调群（作为代数）在各自相应的
  Novikov 域 $\Lambda^{\mathbf{k},\Gamma_\omega}$ 上都是半单的。
\end{wexer}
% [Tutor Note: 含糊点拨——练习 8.4.1 是本节代数性质的第一个应用：半单性（等价于 $\mathcal{E}$
%  可逆）在 8.5 节之后会与「量子上同调不含幂零元」这一几何现象相联系（从而给出辛流形本身的
%  某种刚性）。验证时注意：(a) 对 $\CP^n$，$\mathrm{QH}^*$ 是域 $\mathbf{k}[s,T]/\langle s^{n+1}=T\rangle$
%  的分式化，故半单性可直接验证；(b) 对 $S^2\times S^2$，需先算出 $a,b$ 的对偶
%  $a^{\#},b^{\#}$（用 (8.35)），再检验 $\mathcal{E}=1*1^{\#}+a*a^{\#}+b*b^{\#}+(a\cup b)*(a\cup b)^{\#}$
%  的可逆性。原书把该验证留作练习，读者应至少核对 $\CP^1$ 与 $\CP^n$ 两种情形。]
